			\TemplateSection[K0E910E990E8A0E320E700E25010112010100]{String}
				\TemplateSubsection[KD1E62BCF1E0DC0C60DCCD91CC3A3130101010100]{最小表示法}
					两种不同方法实现的最小表示法
					\TemplateCode{cpp}{src/String/最小表示法.cpp}
				\TemplateSubsection[K0E510E020E700E020E0A0E2C0E210E8A010112010100]{Manacher}
					\TemplateCode{cpp}{src/String/Manacher.cpp}
				\TemplateSubsection[K0E360E510E7E072F07020E210EA60E360E510E7E010112121202020202121212010100]{KMP, exKMP}
					\TemplateCode{cpp}{src/String/KMP.cpp}
				\TemplateSubsection[K0E090E7C0E8A0E1A0E210E8ACCF440010112010100]{Border树}
					实际上还是kmp构建的失配树(Border树),加上了slink指针等附加结构\\
					slink的含义与PAM+dp 中的slink相似,都是指向公差diff与自身不等的上一个border。\\
					slink将原串的border结构划分成 $\log n$ 段等差数列,diff是公差。
					\TemplateCode{cpp}{src/String/BorderTree.cpp}
				% \subsection{exKMP (zjj)}
				% 	\inputminted{cpp}{src/zjj/exkmp.cpp}
				\TemplateSubsection[K0E020E0A0702D1B510C32E0DC5AF3101011212010100]{AC 自动机}
					\TemplateCode{cpp}{src/String/AC 自动机.cpp}
					\TemplateCode{cpp}{src/String/acamoi.cpp}
				\TemplateSubsection[K0E480EA70E700E1A0E7C0E7007020EA40E7C0E8A0E1A07020E1A0E210E0A0E7C0E510E7E0E7C0E910E320E990E320E7C0E70010112020202020202120202020212010100]{Lyndon Word Decomposition}
					\TemplateCode{cpp}{src/String/Lyndon Word.cpp}
					\TemplateCode{cpp}{src/String/Lyndon_CFL.cpp}
				\TemplateSubsection[KC5460DD1961CD1B510C32E0DC5AF310101010100]{后缀自动机}
					\TemplateCode{cpp}{src/String/SAM.cpp}
				\TemplateSubsection[KC5460DD1961CCCF4C4D1D1190101010100]{后缀数组}
					\TemplateCode{cpp}{src/String/SA.cpp}
				\TemplateSubsection[K0E910E9F0E230E230E320EA607020E090E020E480E020E700E0A0E210E1A07020E990E8A0E210E210702C5460DD1961CCAD113C52A70CCF44001011202020202020212020202020202020212010100]{Suffix Balanced Tree 后缀平衡树}
					\TemplateCode{cpp}{src/String/后缀平衡树.cpp}
				%\subsection{Total LCS 子串对目标求 LCS}
				%\inputminted{cpp}{src/String/Total LCS.cpp}
				\TemplateSubsection[KC4AE0DCFF41CC5460DD1961CD1B510C32E0DC5AF310727C7752ECF062B072A0101010100]{广义后缀自动机(离线)}
					使用insert方法插入。全部插入完后需要build。
					\TemplateCode{cpp}{src/String/GSAM.cpp}
				\TemplateSubsection[KC4AE0DCFF41CD0861CCF062B07020E910E020E5101010202020202121212010100]{广义在线 SAM}
					\TemplateCode{cpp}{src/String/generalizedSAM.cpp}
				\TemplateSubsection[KC58813CEB110D1B510C32E0DC5AF310101010100]{回文自动机}
					1 的子树是长为偶数的串, 0 的子树是长为奇数的. 1 代表空串, 0 没有意义.
					\TemplateCode{cpp}{src/String/PAM.cpp}
					\TemplateCode{cpp}{src/String/jianglyPAM.cpp}
				\TemplateSubsection[KC58813CEB110D1B510C32E0DC5AF310702080307020E1A0E7E0101010100]{回文自动机 + dp}
					使用slink指针优化dp.
					\TemplateCode{cpp}{src/String/PAM_slink.cpp}
				\TemplateSubsection[K0E8A0E9F0E700E91010112010100]{Runs}
					\TemplateCode{cpp}{src/String/Runs.cpp}
					学习command\_block 的blog的runs
					\TemplateCode{cpp}{src/String/Runs_hash.cpp}
				%\subsection{双端插入回文树 (zjj)}
				%\inputminted{cpp}{src/String/PAM_bidirectional.cpp}
				\TemplateSubsection[KD1B516C3FF52C2111307020E2C0E020E910E2C01010202020212010100]{字符串 Hash}
				%\texttt{\input{|"python src/Miscellany/randomprimes.py 3e2 1e6 1e12 1e13 2e13 1e15 1e17 5e17 1e18 2e18"}}
				\TemplateCode{cpp}{src/String/hash.cpp}	
				使用二分hash求lcp/lcs
				\TemplateCode{cpp}{src/String/Hash_lcplcs.cpp}	
				\TemplateSubsection[K0E910E990E8A0E320E700E2507020E0A0E7C0E700E0A0E480E9F0E910E320E7C0E700E9101011202020202020212010100]{String Conclusions}
					\TemplateNote{src/tbr/string-Conclusions.tex}
			\newpage
